% =============================================================================
%  progress-report-02.tex — รายงานความก้าวหน้าโครงการ ULMs ครั้งที่ 2
%  *** ฉบับนี้คือฉบับที่ส่งอาจารย์จริง ***
%
%  build:  cd ~/Projects/SoftwareEn/docs && latexmk progress-report-02.tex
%          ผลลัพธ์ออกที่ build/progress-report-02.pdf
%
%  หลักการเขียนฉบับนี้ — ยกมาจาก progress-report-01.tex อ่านก่อนแก้
%    1) สั้น อาจารย์มีเวลาจำกัด ไม่มีสารบัญ
%    2) มี "แผนปัจจุบัน" เส้นเดียวเท่านั้น ห้ามพล็อตแผนก่อนปรับปรุงลงในกราฟ
%    3) อ้างถึงได้เฉพาะ "รายงานความก้าวหน้าครั้งที่ 1 (27 ส.ค. 2569)" ซึ่งส่ง
%       อาจารย์ไปแล้วจริง — ห้ามอ้างฉบับร่างภายใน (progress.tex / progress2.tex
%       / progress3.tex) ซึ่งไม่เคยส่ง
%    4) รายงานความคืบหน้าที่เกิดขึ้นแล้ว ไม่มีข้อเสนอขอตัดขอบเขต
%    5) ตัดหัวข้อที่ไม่จำเป็นต่อการประเมินความคืบหน้าออก ได้แก่ ปัญหาอุปสรรค
%       การประเมินความเสี่ยง และทรัพยากรที่ต้องการเพิ่ม ข้อจำกัดที่ตรวจพบ
%       ยังคงอยู่ แต่ย้ายไปอยู่ใต้ผลงานแต่ละฝั่ง
%    6) ห้ามรายงานผลงานเป็นจำนวนบรรทัดโค้ดหรือจำนวนหน้าเอกสาร ตัววัดพวกนี้
%       ไม่สะท้อนความก้าวหน้าจริง ให้บอกว่า "ทำอะไรได้แล้ว" แทน
%    7) อุปกรณ์ระดับ T3 (ของติดที่ / การจองสถานที่) อยู่ในขอบเขตที่ต้องทำ
%       ไม่ใช่งานที่ตัดออก ต้องปรากฏในแผน Sprint
%    8) เนื้อความใช้ตัวอักษรน้ำหนักเดียวกันทั้งหมด ห้ามใช้ \textbf เน้นคำกลาง
%       ย่อหน้า ตัวหนาสงวนไว้ให้หัวข้อ หัวตาราง และหน้าปกเท่านั้น
%    9) ศัพท์ที่วงการพัฒนามีชื่อเรียกอยู่แล้ว ใช้คำเดิม ห้ามแปลเป็นไทย
%       เปลี่ยนทั้งฉบับเมื่อ 4 ก.ย. 2569 ตามที่ SA สั่ง เพราะคำแปลอ่านแล้ว
%       เดาไม่ออกว่าหมายถึงอะไร เช่น "โต๊ะพิจารณาอนุมัติ" ที่แปลจาก approval
%       queue หรือ "ส่วนติดต่อโปรแกรม" ที่แปลจาก API
%
%       คำที่ใช้เป็นอังกฤษ — contract · procedure · domain · main · branch ·
%       repository · CI · migrate · build · typecheck · test · test runner ·
%       mock · seed · session · audit log · cron job · unique constraint ·
%       buffer · router · service · store · API · backend · frontend ·
%       notification · AC check · smoke test · database schema · approval queue
%
%       *** เว้นวรรคหน้าและหลังทุกครั้งที่ชนอักษรไทย *** ภาษาไทยไม่มีช่องไฟ
%       ระหว่างคำ ถ้าไม่เว้นจะกลายเป็น "รวมเข้าmain" ซึ่งอ่านไม่ออก
%       และห้ามรัน regex ยุบช่องไฟ (" {2,}" -> " ") ทับทั้งไฟล์ เพราะจะกิน
%       การเยื้องของคอมเมนต์บล็อกนี้ไปด้วย — พลาดมาแล้วสองรอบ
%
%       ที่ยังเป็นไทย คือคำของธุรกิจ ไม่ใช่ของการพัฒนา — ผู้ยืม · เจ้าหน้าที่ ·
%       ผู้อนุมัติ · วงจรยืม–คืน · คำขอยืม · การจองห้อง · ระดับอุปกรณ์ T0–T3
%
%  ตัวเลขทุกตัวอ้างสถานะของ origin/main ที่ commit 99da590 ณ 4 ก.ย. 2569 เวลา 02:00
%  ไม่นับโฟลเดอร์ spike-option-b/ และไม่นับงานบน branch ที่ยังไม่ merge
%
%  รอบ 4 ก.ย. — วันรายงานคือวันที่ 8 จาก 14 วันของ Sprint 3 (28 ส.ค. – 10 ก.ย.)
%  ค่าคอลัมน์ "แผน" จึงสอดแทรกเชิงเส้นระหว่างแถว S2 กับ S3 ในตารางภาคผนวก
%  ที่สัดส่วน 8/14 ไม่ได้อ่านตรงจากแถวใดแถวหนึ่งเหมือนรอบก่อน
%
%  เตือนก่อนแก้รอบหน้า: ทีมพุชงานเข้า main ดึกมาก (รอบนี้มี 11 commit เข้ามา
%  ระหว่าง 3 ก.ย. 20:48 ถึง 4 ก.ย. 01:46 ซึ่งเปลี่ยนจำนวนหน้าที่ต่อข้อมูลจริง
%  จาก 6 เป็น 13) ต้อง git fetch แล้วนับใหม่ทันทีก่อนสร้าง PDF ฉบับที่จะส่ง
%
%  ที่มาของตัวเลข ตรวจซ้ำได้จาก repo:
%    78 procedure  = ผลรวม @Query/@Mutation ใน backend/src/*/*.router.ts 9 ไฟล์
%    227 test      = เลขที่ jest รายงานจาก `npm test` (ไม่ใช่การนับบล็อก it()
%                    ในไฟล์ ซึ่งได้ 202) ในเครื่องผ่าน 201 ตก 26 เพราะไม่มี
%                    Postgres — บน CI ผ่านหมด
%    44 AC check   = `npm run test:ac` ฝั่ง frontend รันแล้วผ่านทั้งหมด
%    24 จาก 30 หน้า = frontend/src/features/**/*-page.tsx ที่ไม่ใช้ PlaceholderPage
%    13 หน้าอ่านข้อมูลจริง = หน้าที่เรียก tRPC เอง หรือผ่าน hook ที่เรียก useTRPC
%    15 procedure ที่เรียกไม่ได้ = FE ประกาศ 71 (11 domain) เทียบ BE 78 (9 domain)
% =============================================================================
% =============================================================================
%  ulms-preamble.tex — preamble ที่ใช้ร่วมกันทุกเอกสารของโครงการ ULMs
%
%  ใช้โดย:  main.tex (ข้อเสนอโครงการ) · progress.tex (รายงานความก้าวหน้า)
%
%  เอกสารที่เรียกใช้ต้องทำสองอย่างหลัง % =============================================================================
%  ulms-preamble.tex — preamble ที่ใช้ร่วมกันทุกเอกสารของโครงการ ULMs
%
%  ใช้โดย:  main.tex (ข้อเสนอโครงการ) · progress.tex (รายงานความก้าวหน้า)
%
%  เอกสารที่เรียกใช้ต้องทำสองอย่างหลัง % =============================================================================
%  ulms-preamble.tex — preamble ที่ใช้ร่วมกันทุกเอกสารของโครงการ ULMs
%
%  ใช้โดย:  main.tex (ข้อเสนอโครงการ) · progress.tex (รายงานความก้าวหน้า)
%
%  เอกสารที่เรียกใช้ต้องทำสองอย่างหลัง \input{ulms-preamble}:
%    1) \hypersetup{pdftitle={...}}  — ชื่อเอกสารใน metadata ของ PDF
%    2) \begin{document}
%  (pdftitle ถูกย้ายออกจากไฟล์นี้ เพราะเป็นค่าเฉพาะของแต่ละเอกสาร)
%
%  แก้ที่นี่ที่เดียว = มีผลกับทุกเอกสาร ห้ามก๊อป preamble ไปวางซ้ำ
% =============================================================================
% =============================================================================
%  main.tex  —  เอกสารข้อเสนอโครงการ (Academic Project Proposal)
%  ระบบบริหารจัดการการยืม–คืนอุปกรณ์มหาวิทยาลัย (ULMs)
%
%  IMPORTANT: This document contains Thai text and MUST be compiled with
%  XeLaTeX (not pdflatex). Thai needs a Unicode font + ICU line-breaking.
%      Build:  latexmk main.tex        (uses .latexmkrc -> xelatex)
%      or:     xelatex main.tex        (run twice for the table of contents)
% =============================================================================

\documentclass[11pt, a4paper]{article}

% ---- Fonts & Thai support ----------------------------------------------------
% fontspec lets XeLaTeX use any system font. Laksaman is a complete serif face
% that covers BOTH Thai and Latin glyphs, so one \setmainfont handles the whole
% (Thai + English) document without font switching.
\usepackage{fontspec}

% Thai has no spaces between words, so TeX cannot find line-break points on its
% own. These two XeTeX primitives switch on ICU's dictionary-based Thai word
% segmentation, which is what makes Thai paragraphs wrap correctly.
\XeTeXlinebreaklocale "th"
\XeTeXlinebreakskip = 0pt plus 0.1pt

\setmainfont{Laksaman}

% ---- Page geometry & paragraph style ----------------------------------------
\usepackage[margin=2.5cm]{geometry}
\usepackage{parskip}          % block paragraphs (no indent, gap between) — common
                              % and more readable for Thai documents
\linespread{1.15}             % a little extra leading; Thai has tall tone marks

% ---- Colours (monochrome / grayscale palette) ------------------------------
\usepackage[table]{xcolor}    % 'table' option enables \rowcolor for table headers
% Monochrome (grayscale) palette. Names kept as-is so every reference still
% works; only the tones changed from the original green/blue accents.
\definecolor{kugreen}{HTML}{1A1A1A}   % near-black — section headings & title
\definecolor{kudark}{HTML}{333333}    % dark gray — subsection/subsubsection titles
\definecolor{headbg}{HTML}{EAEAEA}    % light gray background for table header rows
\definecolor{linkblue}{HTML}{404040}  % medium-dark gray — links / citations

% ---- Section heading styling -------------------------------------------------
% 'nobottomtitles*' stops a heading from being printed near the bottom of a
% page: if less than \bottomtitlespace of room is left, the heading (and its
% following text) moves to the next page. This is titlesec's built-in, robust
% replacement for the earlier \needspace hack — the latter injected vertical
% glue that clashed with longtable page-splitting ("Infinite glue shrinkage").
\usepackage[nobottomtitles*]{titlesec}
\titleformat{\section}
  {\color{kugreen}\Large\bfseries}{\thesection}{0.6em}{}
\titleformat{\subsection}
  {\color{kudark}\large\bfseries}{\thesubsection}{0.5em}{}
\titleformat{\subsubsection}
  {\color{kudark}\normalsize\bfseries}{\thesubsubsection}{0.5em}{}
\titlespacing*{\section}{0pt}{1.4\baselineskip}{0.5\baselineskip}

% Minimum room a heading needs at the bottom of a page; if only ~1–3 lines are
% left (less than this), the heading is bumped to the next page instead of being
% stranded above its table/text.
\renewcommand{\bottomtitlespace}{4\baselineskip}

% ---- Hierarchical indentation by heading level ------------------------------
% Each heading level shifts BOTH the heading and the body text under it further
% to the right: section = 0, subsection (x.y) = 1 step, subsubsection (x.y.z) =
% 2 steps. We do this by having the sectioning commands set \leftskip (a global
% left margin that carries through every following line until the next heading
% resets it). \par first, so the *previous* paragraph is not re-indented.
% \leftskip indents plain paragraphs; \@totalleftmargin makes list environments
% (itemize/enumerate) start from the same indent instead of the page margin.
\makeatletter
\newlength{\lvlindent}
\setlength{\lvlindent}{1.8em}          % one indent step
\let\ulmsSecIndent\section
\renewcommand{\section}{\par\global\leftskip=0pt\global\@totalleftmargin=0pt\relax\ulmsSecIndent}
\let\ulmsSubsecIndent\subsection
\renewcommand{\subsection}{\par\global\leftskip=\lvlindent\global\@totalleftmargin=\lvlindent\relax\ulmsSubsecIndent}
\let\ulmsSubsubsecIndent\subsubsection
\renewcommand{\subsubsection}{\par\global\leftskip=2\lvlindent\global\@totalleftmargin=2\lvlindent\relax\ulmsSubsubsecIndent}
\makeatother

% Indent the headings themselves to match (titlesec ignores \leftskip for the
% heading line, so its left margin is set here via the first \titlespacing arg).
\titlespacing*{\subsection}   {\lvlindent} {1.0\baselineskip}{0.4\baselineskip}
\titlespacing*{\subsubsection}{2\lvlindent}{0.8\baselineskip}{0.3\baselineskip}

% ---- Lists (tighter, custom bullets) ----------------------------------------
\usepackage{enumitem}
\setlist[itemize]{leftmargin=1.6em, itemsep=2pt, topsep=3pt}
\setlist[enumerate]{leftmargin=1.9em, itemsep=2pt, topsep=3pt}

% ---- Tables ------------------------------------------------------------------
\usepackage{array}
\usepackage{booktabs}         % \toprule / \midrule / \bottomrule
\usepackage{tabularx}         % auto-width columns that wrap long Thai text
\usepackage{longtable}        % tables that can break across pages
\usepackage{multirow}         % merged cells for the team-structure table

% Left-aligned wrapping column types:
%   L{width} — fixed-width ragged-right paragraph column
%   Y        — flexible ragged-right X column (for tabularx)
\newcolumntype{L}[1]{>{\raggedright\arraybackslash}p{#1}}
\newcolumntype{Y}{>{\raggedright\arraybackslash}X}
\renewcommand{\arraystretch}{1.3}

% Helpers for a styled header row inside tables.
%   \hrow  — starts a header row and shades it (must be FIRST in the row;
%            colortbl allows only one \rowcolor per row).
%   \thead — bold header-cell text.
\newcommand{\hrow}{\rowcolor{headbg}}
\newcommand{\thead}[1]{\textbf{#1}}

% \begin{keeptable} … \end{keeptable} — keeps a table and the bold label above
% it on the SAME page (used by the data dictionary in §5). A tabularx cannot
% split across pages, so LaTeX is free to break between the label and its
% table, stranding the label alone at the bottom of a page; wrapping the pair
% in a minipage makes them one unbreakable block.
%   * \leftskip is zeroed inside: the minipage box is already placed at the
%     current heading indent, so leaving it set would indent the content twice
%     and push the table past the right margin.
%   * \parskip is restored because minipage's \@parboxrestore zeroes it, which
%     would glue the table straight onto the label line.
% Tables inside must use \linewidth (= the minipage width), not \textwidth.
\newenvironment{keeptable}
  {\par\medskip\noindent
   \begin{minipage}{\dimexpr\textwidth-\leftskip\relax}%
   \leftskip=0pt\relax
   \setlength{\parskip}{0.5\baselineskip plus 2pt}}
  {\end{minipage}\par\medskip}

% ---- Table of contents styling ----------------------------------------------
% By default LaTeX gives dotted leaders to subsections but NOT to sections
% (they print bold with just a right-aligned page number). tocloft lets us turn
% on dots for the section level too, so every TOC line runs "….." to its page.
\usepackage{tocloft}
\renewcommand{\cftsecdotsep}{\cftdotsep}   % dotted leaders for section entries
\renewcommand{\cftsecleader}{\cftdotfill{\cftsecdotsep}}
\renewcommand{\cftsecpagefont}{\normalfont}  % keep page numbers un-bold

% ---- Figures -----------------------------------------------------------------
\usepackage{graphicx}
\graphicspath{{figures/}}

% Landscape (rotated) full-page floats — needed by the ER diagram in §5, which
% is far wider than it is tall and would be unreadable scaled to \textwidth.
\usepackage{rotating}

% ---- Flowcharts / diagrams (TikZ) -------------------------------------------
% Used for the borrow–return workflow diagram (§5). Libraries: geometric shapes
% for decision diamonds, arrows.meta for arrowheads, positioning for relative
% node placement, calc/fit/backgrounds for routing loops and grouping.
\usepackage{tikz}
\usetikzlibrary{shapes.geometric, arrows.meta, positioning, calc, fit, backgrounds}

% ---- Links & bookmarks -------------------------------------------------------
\usepackage{hyperref}
\hypersetup{
  colorlinks=true, linkcolor=black, citecolor=linkblue, urlcolor=linkblue,
  pdflang={th},
}

% Thai labels for auto-generated headings.
\renewcommand{\contentsname}{สารบัญ (Table of Contents)}
\renewcommand{\tablename}{ตาราง}
\renewcommand{\figurename}{รูปที่}

:
%    1) \hypersetup{pdftitle={...}}  — ชื่อเอกสารใน metadata ของ PDF
%    2) \begin{document}
%  (pdftitle ถูกย้ายออกจากไฟล์นี้ เพราะเป็นค่าเฉพาะของแต่ละเอกสาร)
%
%  แก้ที่นี่ที่เดียว = มีผลกับทุกเอกสาร ห้ามก๊อป preamble ไปวางซ้ำ
% =============================================================================
% =============================================================================
%  main.tex  —  เอกสารข้อเสนอโครงการ (Academic Project Proposal)
%  ระบบบริหารจัดการการยืม–คืนอุปกรณ์มหาวิทยาลัย (ULMs)
%
%  IMPORTANT: This document contains Thai text and MUST be compiled with
%  XeLaTeX (not pdflatex). Thai needs a Unicode font + ICU line-breaking.
%      Build:  latexmk main.tex        (uses .latexmkrc -> xelatex)
%      or:     xelatex main.tex        (run twice for the table of contents)
% =============================================================================

\documentclass[11pt, a4paper]{article}

% ---- Fonts & Thai support ----------------------------------------------------
% fontspec lets XeLaTeX use any system font. Laksaman is a complete serif face
% that covers BOTH Thai and Latin glyphs, so one \setmainfont handles the whole
% (Thai + English) document without font switching.
\usepackage{fontspec}

% Thai has no spaces between words, so TeX cannot find line-break points on its
% own. These two XeTeX primitives switch on ICU's dictionary-based Thai word
% segmentation, which is what makes Thai paragraphs wrap correctly.
\XeTeXlinebreaklocale "th"
\XeTeXlinebreakskip = 0pt plus 0.1pt

\setmainfont{Laksaman}

% ---- Page geometry & paragraph style ----------------------------------------
\usepackage[margin=2.5cm]{geometry}
\usepackage{parskip}          % block paragraphs (no indent, gap between) — common
                              % and more readable for Thai documents
\linespread{1.15}             % a little extra leading; Thai has tall tone marks

% ---- Colours (monochrome / grayscale palette) ------------------------------
\usepackage[table]{xcolor}    % 'table' option enables \rowcolor for table headers
% Monochrome (grayscale) palette. Names kept as-is so every reference still
% works; only the tones changed from the original green/blue accents.
\definecolor{kugreen}{HTML}{1A1A1A}   % near-black — section headings & title
\definecolor{kudark}{HTML}{333333}    % dark gray — subsection/subsubsection titles
\definecolor{headbg}{HTML}{EAEAEA}    % light gray background for table header rows
\definecolor{linkblue}{HTML}{404040}  % medium-dark gray — links / citations

% ---- Section heading styling -------------------------------------------------
% 'nobottomtitles*' stops a heading from being printed near the bottom of a
% page: if less than \bottomtitlespace of room is left, the heading (and its
% following text) moves to the next page. This is titlesec's built-in, robust
% replacement for the earlier \needspace hack — the latter injected vertical
% glue that clashed with longtable page-splitting ("Infinite glue shrinkage").
\usepackage[nobottomtitles*]{titlesec}
\titleformat{\section}
  {\color{kugreen}\Large\bfseries}{\thesection}{0.6em}{}
\titleformat{\subsection}
  {\color{kudark}\large\bfseries}{\thesubsection}{0.5em}{}
\titleformat{\subsubsection}
  {\color{kudark}\normalsize\bfseries}{\thesubsubsection}{0.5em}{}
\titlespacing*{\section}{0pt}{1.4\baselineskip}{0.5\baselineskip}

% Minimum room a heading needs at the bottom of a page; if only ~1–3 lines are
% left (less than this), the heading is bumped to the next page instead of being
% stranded above its table/text.
\renewcommand{\bottomtitlespace}{4\baselineskip}

% ---- Hierarchical indentation by heading level ------------------------------
% Each heading level shifts BOTH the heading and the body text under it further
% to the right: section = 0, subsection (x.y) = 1 step, subsubsection (x.y.z) =
% 2 steps. We do this by having the sectioning commands set \leftskip (a global
% left margin that carries through every following line until the next heading
% resets it). \par first, so the *previous* paragraph is not re-indented.
% \leftskip indents plain paragraphs; \@totalleftmargin makes list environments
% (itemize/enumerate) start from the same indent instead of the page margin.
\makeatletter
\newlength{\lvlindent}
\setlength{\lvlindent}{1.8em}          % one indent step
\let\ulmsSecIndent\section
\renewcommand{\section}{\par\global\leftskip=0pt\global\@totalleftmargin=0pt\relax\ulmsSecIndent}
\let\ulmsSubsecIndent\subsection
\renewcommand{\subsection}{\par\global\leftskip=\lvlindent\global\@totalleftmargin=\lvlindent\relax\ulmsSubsecIndent}
\let\ulmsSubsubsecIndent\subsubsection
\renewcommand{\subsubsection}{\par\global\leftskip=2\lvlindent\global\@totalleftmargin=2\lvlindent\relax\ulmsSubsubsecIndent}
\makeatother

% Indent the headings themselves to match (titlesec ignores \leftskip for the
% heading line, so its left margin is set here via the first \titlespacing arg).
\titlespacing*{\subsection}   {\lvlindent} {1.0\baselineskip}{0.4\baselineskip}
\titlespacing*{\subsubsection}{2\lvlindent}{0.8\baselineskip}{0.3\baselineskip}

% ---- Lists (tighter, custom bullets) ----------------------------------------
\usepackage{enumitem}
\setlist[itemize]{leftmargin=1.6em, itemsep=2pt, topsep=3pt}
\setlist[enumerate]{leftmargin=1.9em, itemsep=2pt, topsep=3pt}

% ---- Tables ------------------------------------------------------------------
\usepackage{array}
\usepackage{booktabs}         % \toprule / \midrule / \bottomrule
\usepackage{tabularx}         % auto-width columns that wrap long Thai text
\usepackage{longtable}        % tables that can break across pages
\usepackage{multirow}         % merged cells for the team-structure table

% Left-aligned wrapping column types:
%   L{width} — fixed-width ragged-right paragraph column
%   Y        — flexible ragged-right X column (for tabularx)
\newcolumntype{L}[1]{>{\raggedright\arraybackslash}p{#1}}
\newcolumntype{Y}{>{\raggedright\arraybackslash}X}
\renewcommand{\arraystretch}{1.3}

% Helpers for a styled header row inside tables.
%   \hrow  — starts a header row and shades it (must be FIRST in the row;
%            colortbl allows only one \rowcolor per row).
%   \thead — bold header-cell text.
\newcommand{\hrow}{\rowcolor{headbg}}
\newcommand{\thead}[1]{\textbf{#1}}

% \begin{keeptable} … \end{keeptable} — keeps a table and the bold label above
% it on the SAME page (used by the data dictionary in §5). A tabularx cannot
% split across pages, so LaTeX is free to break between the label and its
% table, stranding the label alone at the bottom of a page; wrapping the pair
% in a minipage makes them one unbreakable block.
%   * \leftskip is zeroed inside: the minipage box is already placed at the
%     current heading indent, so leaving it set would indent the content twice
%     and push the table past the right margin.
%   * \parskip is restored because minipage's \@parboxrestore zeroes it, which
%     would glue the table straight onto the label line.
% Tables inside must use \linewidth (= the minipage width), not \textwidth.
\newenvironment{keeptable}
  {\par\medskip\noindent
   \begin{minipage}{\dimexpr\textwidth-\leftskip\relax}%
   \leftskip=0pt\relax
   \setlength{\parskip}{0.5\baselineskip plus 2pt}}
  {\end{minipage}\par\medskip}

% ---- Table of contents styling ----------------------------------------------
% By default LaTeX gives dotted leaders to subsections but NOT to sections
% (they print bold with just a right-aligned page number). tocloft lets us turn
% on dots for the section level too, so every TOC line runs "….." to its page.
\usepackage{tocloft}
\renewcommand{\cftsecdotsep}{\cftdotsep}   % dotted leaders for section entries
\renewcommand{\cftsecleader}{\cftdotfill{\cftsecdotsep}}
\renewcommand{\cftsecpagefont}{\normalfont}  % keep page numbers un-bold

% ---- Figures -----------------------------------------------------------------
\usepackage{graphicx}
\graphicspath{{figures/}}

% Landscape (rotated) full-page floats — needed by the ER diagram in §5, which
% is far wider than it is tall and would be unreadable scaled to \textwidth.
\usepackage{rotating}

% ---- Flowcharts / diagrams (TikZ) -------------------------------------------
% Used for the borrow–return workflow diagram (§5). Libraries: geometric shapes
% for decision diamonds, arrows.meta for arrowheads, positioning for relative
% node placement, calc/fit/backgrounds for routing loops and grouping.
\usepackage{tikz}
\usetikzlibrary{shapes.geometric, arrows.meta, positioning, calc, fit, backgrounds}

% ---- Links & bookmarks -------------------------------------------------------
\usepackage{hyperref}
\hypersetup{
  colorlinks=true, linkcolor=black, citecolor=linkblue, urlcolor=linkblue,
  pdflang={th},
}

% Thai labels for auto-generated headings.
\renewcommand{\contentsname}{สารบัญ (Table of Contents)}
\renewcommand{\tablename}{ตาราง}
\renewcommand{\figurename}{รูปที่}

:
%    1) \hypersetup{pdftitle={...}}  — ชื่อเอกสารใน metadata ของ PDF
%    2) \begin{document}
%  (pdftitle ถูกย้ายออกจากไฟล์นี้ เพราะเป็นค่าเฉพาะของแต่ละเอกสาร)
%
%  แก้ที่นี่ที่เดียว = มีผลกับทุกเอกสาร ห้ามก๊อป preamble ไปวางซ้ำ
% =============================================================================
% =============================================================================
%  main.tex  —  เอกสารข้อเสนอโครงการ (Academic Project Proposal)
%  ระบบบริหารจัดการการยืม–คืนอุปกรณ์มหาวิทยาลัย (ULMs)
%
%  IMPORTANT: This document contains Thai text and MUST be compiled with
%  XeLaTeX (not pdflatex). Thai needs a Unicode font + ICU line-breaking.
%      Build:  latexmk main.tex        (uses .latexmkrc -> xelatex)
%      or:     xelatex main.tex        (run twice for the table of contents)
% =============================================================================

\documentclass[11pt, a4paper]{article}

% ---- Fonts & Thai support ----------------------------------------------------
% fontspec lets XeLaTeX use any system font. Laksaman is a complete serif face
% that covers BOTH Thai and Latin glyphs, so one \setmainfont handles the whole
% (Thai + English) document without font switching.
\usepackage{fontspec}

% Thai has no spaces between words, so TeX cannot find line-break points on its
% own. These two XeTeX primitives switch on ICU's dictionary-based Thai word
% segmentation, which is what makes Thai paragraphs wrap correctly.
\XeTeXlinebreaklocale "th"
\XeTeXlinebreakskip = 0pt plus 0.1pt

\setmainfont{Laksaman}

% ---- Page geometry & paragraph style ----------------------------------------
\usepackage[margin=2.5cm]{geometry}
\usepackage{parskip}          % block paragraphs (no indent, gap between) — common
                              % and more readable for Thai documents
\linespread{1.15}             % a little extra leading; Thai has tall tone marks

% ---- Colours (monochrome / grayscale palette) ------------------------------
\usepackage[table]{xcolor}    % 'table' option enables \rowcolor for table headers
% Monochrome (grayscale) palette. Names kept as-is so every reference still
% works; only the tones changed from the original green/blue accents.
\definecolor{kugreen}{HTML}{1A1A1A}   % near-black — section headings & title
\definecolor{kudark}{HTML}{333333}    % dark gray — subsection/subsubsection titles
\definecolor{headbg}{HTML}{EAEAEA}    % light gray background for table header rows
\definecolor{linkblue}{HTML}{404040}  % medium-dark gray — links / citations

% ---- Section heading styling -------------------------------------------------
% 'nobottomtitles*' stops a heading from being printed near the bottom of a
% page: if less than \bottomtitlespace of room is left, the heading (and its
% following text) moves to the next page. This is titlesec's built-in, robust
% replacement for the earlier \needspace hack — the latter injected vertical
% glue that clashed with longtable page-splitting ("Infinite glue shrinkage").
\usepackage[nobottomtitles*]{titlesec}
\titleformat{\section}
  {\color{kugreen}\Large\bfseries}{\thesection}{0.6em}{}
\titleformat{\subsection}
  {\color{kudark}\large\bfseries}{\thesubsection}{0.5em}{}
\titleformat{\subsubsection}
  {\color{kudark}\normalsize\bfseries}{\thesubsubsection}{0.5em}{}
\titlespacing*{\section}{0pt}{1.4\baselineskip}{0.5\baselineskip}

% Minimum room a heading needs at the bottom of a page; if only ~1–3 lines are
% left (less than this), the heading is bumped to the next page instead of being
% stranded above its table/text.
\renewcommand{\bottomtitlespace}{4\baselineskip}

% ---- Hierarchical indentation by heading level ------------------------------
% Each heading level shifts BOTH the heading and the body text under it further
% to the right: section = 0, subsection (x.y) = 1 step, subsubsection (x.y.z) =
% 2 steps. We do this by having the sectioning commands set \leftskip (a global
% left margin that carries through every following line until the next heading
% resets it). \par first, so the *previous* paragraph is not re-indented.
% \leftskip indents plain paragraphs; \@totalleftmargin makes list environments
% (itemize/enumerate) start from the same indent instead of the page margin.
\makeatletter
\newlength{\lvlindent}
\setlength{\lvlindent}{1.8em}          % one indent step
\let\ulmsSecIndent\section
\renewcommand{\section}{\par\global\leftskip=0pt\global\@totalleftmargin=0pt\relax\ulmsSecIndent}
\let\ulmsSubsecIndent\subsection
\renewcommand{\subsection}{\par\global\leftskip=\lvlindent\global\@totalleftmargin=\lvlindent\relax\ulmsSubsecIndent}
\let\ulmsSubsubsecIndent\subsubsection
\renewcommand{\subsubsection}{\par\global\leftskip=2\lvlindent\global\@totalleftmargin=2\lvlindent\relax\ulmsSubsubsecIndent}
\makeatother

% Indent the headings themselves to match (titlesec ignores \leftskip for the
% heading line, so its left margin is set here via the first \titlespacing arg).
\titlespacing*{\subsection}   {\lvlindent} {1.0\baselineskip}{0.4\baselineskip}
\titlespacing*{\subsubsection}{2\lvlindent}{0.8\baselineskip}{0.3\baselineskip}

% ---- Lists (tighter, custom bullets) ----------------------------------------
\usepackage{enumitem}
\setlist[itemize]{leftmargin=1.6em, itemsep=2pt, topsep=3pt}
\setlist[enumerate]{leftmargin=1.9em, itemsep=2pt, topsep=3pt}

% ---- Tables ------------------------------------------------------------------
\usepackage{array}
\usepackage{booktabs}         % \toprule / \midrule / \bottomrule
\usepackage{tabularx}         % auto-width columns that wrap long Thai text
\usepackage{longtable}        % tables that can break across pages
\usepackage{multirow}         % merged cells for the team-structure table

% Left-aligned wrapping column types:
%   L{width} — fixed-width ragged-right paragraph column
%   Y        — flexible ragged-right X column (for tabularx)
\newcolumntype{L}[1]{>{\raggedright\arraybackslash}p{#1}}
\newcolumntype{Y}{>{\raggedright\arraybackslash}X}
\renewcommand{\arraystretch}{1.3}

% Helpers for a styled header row inside tables.
%   \hrow  — starts a header row and shades it (must be FIRST in the row;
%            colortbl allows only one \rowcolor per row).
%   \thead — bold header-cell text.
\newcommand{\hrow}{\rowcolor{headbg}}
\newcommand{\thead}[1]{\textbf{#1}}

% \begin{keeptable} … \end{keeptable} — keeps a table and the bold label above
% it on the SAME page (used by the data dictionary in §5). A tabularx cannot
% split across pages, so LaTeX is free to break between the label and its
% table, stranding the label alone at the bottom of a page; wrapping the pair
% in a minipage makes them one unbreakable block.
%   * \leftskip is zeroed inside: the minipage box is already placed at the
%     current heading indent, so leaving it set would indent the content twice
%     and push the table past the right margin.
%   * \parskip is restored because minipage's \@parboxrestore zeroes it, which
%     would glue the table straight onto the label line.
% Tables inside must use \linewidth (= the minipage width), not \textwidth.
\newenvironment{keeptable}
  {\par\medskip\noindent
   \begin{minipage}{\dimexpr\textwidth-\leftskip\relax}%
   \leftskip=0pt\relax
   \setlength{\parskip}{0.5\baselineskip plus 2pt}}
  {\end{minipage}\par\medskip}

% ---- Table of contents styling ----------------------------------------------
% By default LaTeX gives dotted leaders to subsections but NOT to sections
% (they print bold with just a right-aligned page number). tocloft lets us turn
% on dots for the section level too, so every TOC line runs "….." to its page.
\usepackage{tocloft}
\renewcommand{\cftsecdotsep}{\cftdotsep}   % dotted leaders for section entries
\renewcommand{\cftsecleader}{\cftdotfill{\cftsecdotsep}}
\renewcommand{\cftsecpagefont}{\normalfont}  % keep page numbers un-bold

% ---- Figures -----------------------------------------------------------------
\usepackage{graphicx}
\graphicspath{{figures/}}

% Landscape (rotated) full-page floats — needed by the ER diagram in §5, which
% is far wider than it is tall and would be unreadable scaled to \textwidth.
\usepackage{rotating}

% ---- Flowcharts / diagrams (TikZ) -------------------------------------------
% Used for the borrow–return workflow diagram (§5). Libraries: geometric shapes
% for decision diamonds, arrows.meta for arrowheads, positioning for relative
% node placement, calc/fit/backgrounds for routing loops and grouping.
\usepackage{tikz}
\usetikzlibrary{shapes.geometric, arrows.meta, positioning, calc, fit, backgrounds}

% ---- Links & bookmarks -------------------------------------------------------
\usepackage{hyperref}
\hypersetup{
  colorlinks=true, linkcolor=black, citecolor=linkblue, urlcolor=linkblue,
  pdflang={th},
}

% Thai labels for auto-generated headings.
\renewcommand{\contentsname}{สารบัญ (Table of Contents)}
\renewcommand{\tablename}{ตาราง}
\renewcommand{\figurename}{รูปที่}



\hypersetup{pdftitle={รายงานความก้าวหน้าโครงการ ULMs ครั้งที่ 2}}

% ---- ย่อหน้า ------------------------------------------------------------------
% preamble โหลด parskip ไว้ ซึ่งบังคับ \parindent=0 แล้วเว้นบรรทัดคั่นแทน
% รายงานฉบับนี้ต้องการย่อหน้าแบบไทย คือเยื้องบรรทัดแรก จึงตั้ง \parindent คืน
\setlength{\parindent}{1.5em}
\setlength{\parskip}{0.35\baselineskip plus 2pt}

% preamble ประกาศ \titlespacing* แบบมีดอกจัน ซึ่ง titlesec ตีความว่าห้ามเยื้อง
% ย่อหน้าแรกที่ตามหลังหัวข้อ เอกสารไทยเยื้องทุกย่อหน้า จึงประกาศซ้ำแบบไม่มี
% ดอกจันด้วยค่าเดิมทุกตัว — ค่าที่ใช้ต้องตรงกับ ulms-preamble.tex เสมอ
\titlespacing{\section} {0pt} {1.4\baselineskip}{0.5\baselineskip}
\titlespacing{\subsection} {\lvlindent} {1.0\baselineskip}{0.4\baselineskip}
\titlespacing{\subsubsection}{2\lvlindent}{0.8\baselineskip}{0.3\baselineskip}

\setlist[itemize]{leftmargin=1.6em, itemsep=2pt, topsep=0.55\baselineskip}
\setlist[enumerate]{leftmargin=1.9em, itemsep=2pt, topsep=0.55\baselineskip}

% ---- การจัดบรรทัด ------------------------------------------------------------
% ภาษาไทยไม่มีการตัดคำด้วยยัติภังค์ และรายงานมีชื่อไฟล์ใน \texttt{} จำนวนมาก
% ซึ่ง TeX ตัดกลางคำไม่ได้ \emergencystretch สั่งให้ยอมยืดช่องไฟแทนการปล่อยล้น
\setlength{\emergencystretch}{3em}

% \zb — จุดตัดบรรทัดกว้างศูนย์ สำหรับแทรกในชื่อยาว ๆ ที่ \texttt ตัดเองไม่ได้
\newcommand{\zb}{\hspace{0pt}}

% ---- การวางรูป --------------------------------------------------------------
% กราฟทั้งสองรูปต้องอยู่ตรงตำแหน่งที่เขียนไว้เป๊ะ ๆ ห้ามลอยข้ามหัวข้อไป
\usepackage{float}

% ---- ตัวแปรของรอบรายงาน ------------------------------------------------------
\newcommand{\reportdate}{4 กันยายน 2569}
\newcommand{\reportperiod}{28 สิงหาคม – 4 กันยายน 2569}
\newcommand{\actualpct}{70\%}
\newcommand{\planpct}{72\%}
\newcommand{\prevactualpct}{52\%}

\begin{document}

% =============================================================================
% TITLE PAGE
% =============================================================================
\begin{titlepage}
\centering
\vspace*{2.8cm}

{\Large\bfseries\color{kugreen} รายงานความก้าวหน้าโครงการ ครั้งที่ 2\par}
\vspace{0.6cm}
{\LARGE\bfseries\color{kugreen} ระบบบริหารจัดการการยืม–คืนอุปกรณ์มหาวิทยาลัย\par}
\vspace{0.35cm}
{\Large\color{kudark} University Lending Management System (ULMs)\par}

\vspace{1.8cm}
\rule{0.72\textwidth}{0.8pt}
\vspace{0.9cm}

\begin{tabular}{@{}r@{\hspace{1.2em}}l@{}}
\textbf{รอบรายงาน} & \reportperiod \\[0.4em]
\textbf{วันที่รายงาน} & \reportdate \\[0.4em]
\textbf{ขอบเขตการนับ} & งานพัฒนาระบบ Sprint 0–5 (ไม่รวมงานก่อน 24 ก.ค.) \\[0.4em]
\textbf{สถานะ Sprint} & วันที่ 8 จาก 14 ของ Sprint 3 (Sprint ที่ 4 จาก 6) \\[0.4em]
\textbf{ความก้าวหน้า} & \actualpct\ เทียบแผน \planpct\ \ (ต่ำกว่าแผน 2\%) \\[0.4em]
\textbf{งบประมาณ} & ใช้ไป 0 บาท จากวงเงินประเมิน 5{,}000 บาท \\[0.4em]
\textbf{กำหนดเสร็จ} & 1 ตุลาคม 2569 (ยังไม่มีการเลื่อน) \\
\end{tabular}

\vspace{0.9cm}
\rule{0.72\textwidth}{0.8pt}

\vfill
{\large\color{kudark} คณะทำงาน ``miro ปั่น 14 บาท''\par}
\vspace{0.3cm}
{\color{kudark} ภาควิชาวิศวกรรมคอมพิวเตอร์ คณะวิศวกรรมศาสตร์\par}
{\color{kudark} มหาวิทยาลัยเกษตรศาสตร์\par}
\vspace{1.2cm}
\end{titlepage}

% =============================================================================
\section{บทสรุป}

รอบรายงานนี้เป็นสัปดาห์แรกของ Sprint 3 และเป็นรอบที่ระยะห่างระหว่างแผนกับผลจริงแคบลงมากที่สุดนับจากเริ่มโครงการ จากต่ำกว่าแผน 13\% ในรอบก่อนเหลือต่ำกว่าแผน 2\% ในรอบนี้ คณะทำงานมีกำลังเต็มตลอดทั้งสัปดาห์และมีงานเข้า main ทุกวัน

ความก้าวหน้าของรอบนี้มาจากสี่เรื่อง เรื่องแรกคืองานที่รอบก่อนรายงานว่า ``เขียนเสร็จแล้วแต่ยังไม่รวมเข้า main'' ถูกรวมเข้า main ครบทุกก้อน ทั้งรายการอุปกรณ์และการค้นหา ระบบเครดิต วงจรยืม–คืนฝั่งเจ้าหน้าที่ การตรวจสภาพ และการอัปโหลดรูปภาพ พร้อมกับตัดสินให้เหลือวิธีจัดระดับอุปกรณ์เพียงวิธีเดียวตามที่ค้างไว้

เรื่องที่สองคือชิ้นส่วนที่หายไปของวงจรยืม–คืนถูกเขียนขึ้นใหม่และเข้า main แล้ว คือ procedure สร้างคำขอยืมของฝั่งผู้ยืม กับ approval queue ของผู้อนุมัติ รอบก่อนรายงานว่าสองสิ่งนี้ไม่มีอยู่บน branch ใดเลย เมื่อรวมกับระบบ notification ที่เข้า main ในวันสุดท้ายของรอบ วงจรยืม–คืนจึงมีโค้ดครบทุกขั้นแล้ว ตั้งแต่ผู้ยืมสร้างคำขอ ผู้อนุมัติพิจารณา เจ้าหน้าที่จ่ายของ ไปจนถึงรับคืนและตรวจสภาพ โดยยังไม่มีการทดสอบใดเดินครบเส้นทางนี้

เรื่องที่สามคือเอกสารข้อกำหนดความต้องการของซอฟต์แวร์ (SRS) และเอกสารการออกแบบระบบ (SDS) ฉบับทางการเสร็จแล้วทั้งสองฉบับ ซึ่งเป็นข้อจำกัดที่รอบก่อนระบุไว้ว่ายังไม่มี มิติเอกสารจึงเป็นมิติเดียวที่เดินนำแผนในรอบนี้

เรื่องที่สี่เกิดขึ้นในสามสิบชั่วโมงสุดท้ายของรอบ คือหน้าจอฝั่งเจ้าหน้าที่และฝั่งผู้อนุมัติถูกต่อเข้ากับข้อมูลจริง ได้แก่ คิวงานของเจ้าหน้าที่ คลังอุปกรณ์ การตรวจสภาพ ค่าตั้งของระบบการยืม และคิวพิจารณาอนุมัติ หน้าจอที่มีเนื้อหาจริงจึงเพิ่มจาก 18 เป็น 24 จาก 30 หน้า และหน้าจอที่อ่านข้อมูลจริงเพิ่มจาก 6 เป็น 13 หน้า

ตัวเลขความก้าวหน้าในรายงานนี้วัดปริมาณงานที่เขียนเสร็จ ไม่ใช่ปริมาณงานที่ยืนยันแล้วว่าทำงานถูก ส่วนที่เขียนขึ้นในสองสัปดาห์หลังยังไม่มี test ครอบคลุม รายละเอียดอยู่ในหัวข้อ \ref{sec:dimension} และหัวข้อระดับหลักฐาน

ข้อจำกัดสำคัญที่สุดของรอบนี้คือวงจรยืม–คืนยังเดินครบผ่านหน้าจอจริงไม่ได้ แต่จุดที่ขาดเปลี่ยนไปจากต้นสัปดาห์ ตอนนี้ฝั่งเจ้าหน้าที่และฝั่งผู้อนุมัติทำงานกับข้อมูลจริงได้แล้ว จุดที่ขาดย้ายไปอยู่ต้นทางคือปุ่มส่งคำขอของผู้ยืม ซึ่งยังบันทึกลง store ในเบราว์เซอร์แทนที่จะส่งไปยังฐานข้อมูล และหน้าส่งมอบของเจ้าหน้าที่ซึ่งยังเป็นหน้าว่าง สองจุดนี้คืองานลำดับแรกของสัปดาห์ถัดไป

% =============================================================================
\section{แผนภาพรวมทั้ง 6 Sprint}
\label{sec:overview}

โครงการแบ่งเป็น 6 Sprint บนกรอบเวลา 10 สัปดาห์ ตารางนี้เป็นภาพรวมของทั้งโครงการ รายละเอียดของรอบรายงานนี้อยู่ในหัวข้อ \ref{sec:done} และของ Sprint ที่เหลืออยู่ในหัวข้อ \ref{sec:next}

\noindent\begin{tabularx}{\dimexpr\textwidth-\leftskip\relax}{@{} L{2.9cm} Y L{1.5cm} @{}}
\toprule
\hrow \thead{Sprint} & \thead{ธีมและงานหลัก} & \thead{สถานะ} \\
\midrule
\textbf{Sprint 0} \newline 24 – 30 ก.ค. & วางฐาน — ปิดขอบเขต · ติดตั้งโครงสร้างพื้นฐานทั้งสองฝั่ง · database schema แกนกลาง · ระบบออกแบบและ mock & ผ่าน \\
\textbf{Sprint 1} \newline 31 ก.ค. – 13 ส.ค. & ล็อก contract และเปิดทาง — ล็อก contract API · ระบบเข้าสู่ระบบและออกจากระบบ · CI · เกณฑ์การยอมรับของวงจรหลัก & ผ่าน \\
\textbf{Sprint 2} \newline 14 – 27 ส.ค. & วงจรยืม–คืนของอุปกรณ์เคลื่อนย้ายได้ — รายการอุปกรณ์และจำนวนคงเหลือ · สร้างคำขอ อนุมัติ ส่งมอบ รับคืน · ข้อมูล seed · เลิกใช้ mock ในเส้นทางหลัก & ไม่ผ่านเกณฑ์ \newline (ครึ่งแรกตรงกับสอบกลางภาค ยกงานที่เหลือไป Sprint 3) \\
\textbf{Sprint 3} \newline 28 ส.ค. – 10 ก.ย. & กฎธุรกิจรายระดับและการจองสถานที่ (T3) — อนุมัติตามระดับอุปกรณ์ · ผูกหมายเลขชิ้น · สล็อตเวลาและการจองห้อง · ตรวจสภาพและระบบเครดิต \newline รับงานที่ยกมาจาก Sprint 2 เป็นลำดับแรก & กำลังดำเนินการ \newline (ผ่านครึ่งทาง · งานที่ยกมาจาก Sprint 2 รับครบแล้ว) \\
\textbf{Sprint 4} \newline 11 – 24 ก.ย. & ข้ามหน่วยงานและงานขัดเงา — ตะกร้าข้ามหน่วยงาน · จองล่วงหน้า · ต่ออายุ · หน้าจอที่เหลือ · การค้นหาและการรองรับมือถือ & ตามแผน \\
\textbf{Sprint 5} \newline 25 ก.ย. – 1 ต.ค. & หยุดเพิ่มความสามารถ — ทดสอบถดถอย · นำขึ้นใช้งานจริง · คู่มือผู้ใช้และคู่มือติดตั้ง · ซ้อมนำเสนอ & ตามแผน \\
\bottomrule
\end{tabularx}

% =============================================================================
\section{สถานะโครงการโดยรวม (เงิน · งาน · เวลา)}

\noindent\begin{tabularx}{\dimexpr\textwidth-\leftskip\relax}{@{} L{1.5cm} L{3.0cm} L{2.4cm} Y @{}}
\toprule
\hrow \thead{ด้าน} & \thead{แผน} & \thead{ผลจริง} & \thead{การประเมิน} \\
\midrule
เวลา & ผ่านไป 6.0 จาก 10 สัปดาห์ (60\%) & ผ่านไป 6.0 จาก 10 สัปดาห์ (60\%) & ตามปฏิทิน วันเสร็จยังเป็น 1 ตุลาคม 2569 \\
งาน & \planpct & \actualpct & ต่ำกว่าแผน 2\% แคบลงจาก 13\% ในรอบก่อน เพราะงานที่ค้างบน branch แยกถูกรวมเข้า main ครบ และหน้าจอฝั่งเจ้าหน้าที่ถูกต่อเข้าข้อมูลจริง รายละเอียดรายมิติอยู่ในหัวข้อ \ref{sec:dimension} \\
เงิน & 5{,}000 บาท (วงเงินประเมินทั้งโครงการ) & 0 บาท (0\%) & ยังไม่มีรายจ่าย รายการที่ประเมินไว้คือค่าเช่า Cloud/Hosting และฐานข้อมูลกลาง ซึ่งจะเกิดขึ้นเมื่อนำระบบขึ้นใช้งานจริงใน Sprint 5 \\
\bottomrule
\end{tabularx}

% =============================================================================
\section{ความก้าวหน้าแยกตามมิติของงาน}
\label{sec:dimension}

การแบ่งมิติเป็นไปตามสไลด์ Week 06 คือแยกส่วนออกแบบ ส่วนพัฒนา ส่วนทดสอบ และส่วนเอกสาร ค่าในคอลัมน์ ``จริง'' ประเมินจากผลงานที่ตรวจสอบได้ใน repository และนับเฉพาะงานที่รวมเข้า main แล้วตามมติของทีม

วันรายงานรอบนี้คือวันที่ 8 จาก 14 วันของ Sprint 3 ค่าในคอลัมน์ ``แผน'' จึงสอดแทรกเชิงเส้นระหว่างค่าเป้าหมาย ณ ขอบ Sprint 2 กับขอบ Sprint 3 ในตารางภาคผนวกตามสัดส่วนนั้น ไม่ได้อ่านตรงจากแถวใดแถวหนึ่ง

\noindent\begin{tabularx}{\dimexpr\textwidth-\leftskip\relax}{@{} L{3.2cm} L{1.1cm} L{1.0cm} L{1.0cm} L{1.2cm} Y @{}}
\toprule
\hrow \thead{มิติงาน} & \thead{น้ำหนัก} & \thead{แผน} & \thead{จริง} & \thead{ผลต่าง} & \thead{ข้อเท็จจริงประกอบ} \\
\midrule
Back-End Design \newline (API + Database Model) & 10\% & 98\% & 95\% & $-3$ & contract ครบทั้ง 9 domain อยู่บน main แล้ว database schema เพิ่มตาราง session, audit log, ร่องรอยการอนุมัติ และ notification ยังเหลือกลุ่มอุทธรณ์กับกลุ่มรายงานที่ยังไม่ได้ออกแบบฝั่ง backend \\
Front-End Design \newline (UI/UX + Business Process) & 10\% & 89\% & 91\% & $+2$ & ออกแบบหน้าจอที่มีเนื้อหาจริงแล้ว 24 จาก 30 หน้า เพิ่มจาก 14 หน้าในรอบก่อน โดยเพิ่มฝั่งเจ้าหน้าที่และฝั่งผู้อนุมัติในช่วงท้ายรอบ \\
Back-End Dev & 30\% & 70\% & 69\% & $-1$ & เขียนครบทุกขั้นของวงจรยืม–คืน พร้อมการจองห้องและระบบ notification แต่ยังไม่มีการทดสอบใดเดินวงจรนี้จริง ยังเหลือการอุทธรณ์ รายงานเชิงสถิติ และค่าตั้งเวลาทำการ \\
Front-End Dev & 25\% & 72\% & 69\% & $-3$ & หน้าจอ 13 หน้าอ่านข้อมูลจริงแล้ว เพิ่มจากเส้นทางเดียวในรอบก่อน อีก 11 หน้าที่มี UI ยังอ่าน mock ยังเหลือหน้าว่าง 6 หน้า และปุ่มส่งคำขอของผู้ยืมยังไม่ได้ต่อกับฐานข้อมูล \\
Test (Front-End + Back-End) & 15\% & 45\% & 32\% & $-13$ & test ฝั่ง backend 227 กรณีและ CI เรียกใช้บนฐานข้อมูลจริงแล้ว แต่ 10 จาก 18 service ยังไม่มี test เลย รวมทั้งวงจรยืม–คืน การอนุมัติ และ notification · ฝั่ง frontend ยังไม่มี test runner \\
Documentation & 10\% & 73\% & 85\% & $+12$ & SRS และ SDS ฉบับทางการเสร็จทั้งสองฉบับ พร้อมเอกสารอ้างอิง API ที่สร้างจากโค้ดจริง \\
\midrule
\hrow \thead{รวมถ่วงน้ำหนัก} & \thead{100\%} & \thead{\planpct} & \thead{\actualpct} & \thead{$-2$} & \\
\bottomrule
\end{tabularx}

มิติที่รอบก่อนระบุว่าต้องเฝ้าระวังมีสองมิติ และทั้งสองมิติขยับขึ้นตามที่ประเมินไว้ มิติ Back-End Dev ซึ่งมีน้ำหนักสูงสุดที่ 30\% ขยับจาก 42\% มาอยู่ที่ 69\% จนแทบไม่ต่างจากแผน สาเหตุตรงกับที่รอบก่อนคาดไว้ คืองานส่วนใหญ่เขียนเสร็จอยู่แล้วและรออยู่บน branch แยก การรวมงานเข้า main จึงให้ผลตอบแทนสูงที่สุด ส่วนมิติ Test ขยับจาก 14\% มาอยู่ที่ 32\% เพราะ CI เรียก test ให้ทำงานจริงแล้ว ไม่ใช่เพียงตรวจว่าโค้ดคอมไพล์ผ่าน

มิติที่ต้องเฝ้าระวังในรอบถัดไปคือ Test ซึ่งห่างจากแผนมากที่สุดที่ $-13$ และเป็นมิติเดียวที่ระยะห่างกว้างขึ้นเมื่อตรวจสอบละเอียด ปัญหาไม่ได้อยู่ที่จำนวน แต่อยู่ที่ตำแหน่ง test ที่มีอยู่กระจุกตัวที่การยืนยันตัวตน การเข้ารหัส และการแปลงข้อมูล ซึ่งเป็นส่วนที่เขียนไว้ตั้งแต่ Sprint 1 ส่วนที่เขียนขึ้นใหม่ในสองสัปดาห์หลัง คือวงจรยืม–คืน การอนุมัติ notification การตรวจสภาพ และระบบผู้ดูแลระบบ ยังไม่มี test แม้แต่กรณีเดียว เมื่อนับเป็น service คือ 10 จาก 18 ส่วนที่ไม่มี test เลย

ผลที่ตามมาคือคำว่า ``วงจรยืม–คืนปิดครบ'' ในรายงานนี้หมายถึงเขียนโค้ดครบทุกขั้นและ CI ยืนยันว่าคอมไพล์และ build ผ่าน ไม่ได้หมายความว่ามีอะไรเดินวงจรนั้นตั้งแต่ต้นจนจบแล้ว ณ วันรายงานยังไม่มีทั้ง test และสคริปต์ใดที่ทำเช่นนั้น

ข้อควรระวังในการอ่านตัวเลข \actualpct\ — ตัวเลขนี้วัดปริมาณงานที่เขียนเสร็จ ตามนิยามของมิติที่สไลด์ Week 06 กำหนดไว้ ซึ่งแยกมิติ Dev ที่วัดว่าเขียนอะไรไปแล้ว ออกจากมิติ Test ที่วัดว่ามีอะไรยืนยันว่ามันถูก ถ้าเปลี่ยนไปใช้เกณฑ์เดียวคือ ``มี test อัตโนมัติครอบคลุม'' ตัวเลขจะลดลงเหลือราว 45\% เพราะฝั่ง backend มี 10 จาก 18 service ที่ไม่มี test เลย และฝั่ง frontend ไม่มีหน้าจอสักหน้าเดียวจาก 30 หน้าที่มี test ครอบคลุม คณะทำงานจึงไม่ถือว่า \actualpct\ แปลว่าระบบใช้งานได้แล้ว \actualpct\ ของงานที่วางแผนไว้

มิติ Front-End Design เป็นมิติที่ขยับเร็วที่สุดในรอบนี้จนแซงแผนไป 2 จุด เพราะหน้าจอฝั่งเจ้าหน้าที่และฝั่งผู้อนุมัติเข้ามาในสามสิบชั่วโมงสุดท้าย ส่วน Front-End Dev ยังตามแผนอยู่ 3 จุด เนื่องจากหน้าจอที่ยังอ่าน mock มีอยู่ 11 หน้า และเส้นทางสร้างคำขอของผู้ยืมยังไม่ได้ต่อกับฐานข้อมูล ระยะห่างนี้ปิดได้ด้วยการต่อหน้าจอเข้ากับความสามารถที่มีอยู่แล้ว ไม่ต้องรอความสามารถใหม่

\begin{figure}[H]
\centering
\begin{tikzpicture}[x=0.062cm, y=1cm]
 % ---- กริดแนวตั้งและสเกลแกนนอน ----
 \foreach \x in {0,20,40,60,80,100} {
 \draw[gray!25] (\x,-0.30) -- (\x,6.15);
 \node[below, font=\scriptsize] at (\x,-0.30) {\x};
 }
 \node[below, font=\scriptsize] at (50,-0.68) {\% ความก้าวหน้าของมิตินั้น};
 \draw[thick] (0,-0.30) -- (0,6.15);

 % ---- แท่งรายมิติ: y / ชื่อ / น้ำหนัก / แผน / จริง ----
 \foreach \y/\name/\wt/\plan/\act in {%
 5.6/{Back-End Design}/10/98/95,
 4.6/{Front-End Design}/10/89/91,
 3.6/{Back-End Dev}/30/70/69,
 2.6/{Front-End Dev}/25/72/69,
 1.6/{Test}/15/45/32,
 0.6/{Documentation}/10/73/85} {
 \node[left, xshift=-0.16cm, font=\scriptsize, align=right] at (0,\y)
 {\name\\[-3pt]{\tiny น้ำหนัก \wt\%}};
 \fill[gray!45] (0,{\y+0.05}) rectangle (\plan,{\y+0.30});
 \fill[kugreen] (0,{\y-0.30}) rectangle (\act,{\y-0.05});
 \node[right, font=\tiny, gray!45!black] at (\plan,{\y+0.175}) {\plan};
 \node[right, font=\tiny, kugreen] at (\act,{\y-0.175}) {\act};
 }

 % ---- คำอธิบายสัญลักษณ์ ----
 \fill[gray!45] (0,-1.36) rectangle (7,-1.54);
 \node[right, font=\scriptsize] at (9,-1.45) {แผน};
 \fill[kugreen] (0,-1.76) rectangle (7,-1.94);
 \node[right, font=\scriptsize] at (9,-1.85) {ผลจริง};
\end{tikzpicture}
\vspace{0.6em}
\caption{ความก้าวหน้ารายมิติ ณ \reportdate\ เทียบกับค่าเป้าหมายตามแผน
 ตัวเลขใต้ชื่อมิติคือน้ำหนักที่ใช้ถ่วงเมื่อรวมเป็นค่าเดียว
 แท่งผลจริงนับเฉพาะงานที่รวมเข้า main แล้ว
 มิติที่ห่างจากแผนมากที่สุดคือ Test ที่ $-13$
 ส่วน Documentation กับ Front-End Design เดินนำแผน ที่ $+12$ และ $+2$
 ตามลำดับ}
\label{fig:dims}
\end{figure}

% =============================================================================
\section{กราฟความก้าวหน้าเทียบแผน}

เส้นแผนสร้างจากแผน Sprint โดยตรง ไม่ได้กำหนดค่าด้วยมือ คือกำหนดค่าเป้าหมายรายมิติ ณ วันสิ้นสุดของแต่ละ Sprint จากรายการงานที่ Sprint นั้นรับไว้ แล้วคูณน้ำหนักรวมกันเป็นค่าสะสม ตารางค่าทั้งหมดอยู่ในภาคผนวก

แกนนอนเริ่มที่สัปดาห์ที่ 0 ซึ่งตรงกับ 24 กรกฎาคม 2569 เป็นวันเริ่มลงมือพัฒนา ไม่ใช่วันเริ่มโครงการ งานข้อเสนอและงานออกแบบที่เสร็จก่อนหน้านั้นจึงไม่ปรากฏบนกราฟ

รูปร่างของเส้นผลจริงในรอบนี้บอกสองเรื่อง เรื่องแรกคือความชันของช่วง 27 สิงหาคมถึง 4 กันยายน เป็นความชันที่มากที่สุดของโครงการจนถึงตอนนี้ และมากกว่าความชันของเส้นแผนในช่วงเดียวกัน เรื่องที่สองคือส่วนหนึ่งของความชันนี้เป็นงานที่ทำเสร็จไปแล้วในรอบก่อนแต่เพิ่งถูกนับเมื่อรวมเข้า main ไม่ใช่งานที่เขียนขึ้นใหม่ทั้งหมดภายในแปดวัน คณะทำงานจึงไม่ถือว่าความชันระดับนี้เป็นอัตราที่ทำซ้ำได้ในรอบถัดไป

\begin{figure}[H]
\centering
\begin{tikzpicture}[x=1.20cm, y=0.058cm]
 % ---- กริดแนวนอน ----
 \foreach \y in {0,20,40,60,80,100}
 \draw[gray!22] (0,\y) -- (10.2,\y);

 % ---- เส้นประกำกับขอบ Sprint ----
 \foreach \x in {0.86, 2.86, 4.86, 6.86, 8.86}
 \draw[gray!40, dashed, thin] (\x,0) -- (\x,108);

 % ---- ป้ายชื่อ Sprint ด้านบน ----
 \foreach \x/\lbl in {0.43/{S0}, 1.86/{S1}, 3.86/{S2}, 5.86/{S3}, 7.86/{S4}, 9.36/{S5}}
 \node[font=\scriptsize\bfseries, gray!50!black] at (\x,114) {\lbl};

 % ---- แกน ----
 \draw[->, thick] (0,0) -- (10.6,0);
 \draw[->, thick] (0,0) -- (0,122) node[above, font=\scriptsize] {\% ความก้าวหน้าสะสม};
 \foreach \y in {0,20,40,60,80,100}
 \draw (0,\y) -- (-0.10,\y) node[left, font=\scriptsize] {\y};

 % ---- ขีดและวันที่ตามขอบ Sprint ----
 \foreach \x/\lbl in {0/{24 ก.ค.}, 0.86/{30 ก.ค.}, 2.86/{13 ส.ค.},
 4.86/{27 ส.ค.}, 6.86/{10 ก.ย.}, 8.86/{24 ก.ย.}, 9.86/{1 ต.ค.}}
 {\draw (\x,0) -- (\x,-3);
 \node[rotate=45, anchor=north east, font=\scriptsize] at (\x,-4) {\lbl};}

 % ---- เส้นแนวตั้งของวันรายงาน (4 ก.ย. = สัปดาห์ที่ 6.0 = วันที่ 8 ของ Sprint 3) ----
 \draw[gray!55, dotted, thin] (6.0,0) -- (6.0,108);

 % ---- เส้นแผน — จุดวางที่ขอบ Sprint ----
 \draw[thick, kudark, densely dashed]
 plot coordinates {(0,1) (0.86,12) (2.86,42) (4.86,65) (6.86,78) (8.86,94) (9.86,100)};
 \foreach \p in {(0,1), (0.86,12), (2.86,42), (4.86,65), (6.86,78), (8.86,94), (9.86,100)}
 \fill[kudark] \p circle (1.6pt);

 % ---- เส้นผลจริง — จุดคือวันที่วัดความก้าวหน้า ----
 % ช่วง 2.86 -> 4.29 ราบ เพราะช่วงสอบกลางภาคไม่มีงานเพิ่ม
 % 4.86 -> 6.0 คือรอบรายงานนี้ ชันที่สุดของโครงการ
 \draw[very thick, kugreen]
 plot coordinates {(0,0) (2.0,24) (2.43,32) (2.86,43) (4.29,43) (4.86,52) (6.0,70)};
 \foreach \p in {(0,0), (2.0,24), (2.43,32), (2.86,43), (4.29,43), (4.86,52), (6.0,70)}
 \fill[kugreen] \p circle (2.2pt);

 % ---- ระยะห่างแผน–ผลจริง ณ วันรายงาน ----
 % ค่าแผนที่ใช้คือ 72 ซึ่งเป็นค่าที่ได้จากการปัดรายมิติก่อนถ่วงน้ำหนักในตารางภาคผนวก
 % (ค่าสอดแทรกดิบของคอลัมน์รวมคือ 72.4 ต่างกัน 0.4 ไม่มีผลกับรูป)
 \draw[gray!70, thick, |-|] (6.0,70) -- (6.0,72);
 \node[right, font=\tiny, gray!70!black] at (6.06,71) {$-2$};

 % ---- ป้ายจุดรอบก่อนและจุดปัจจุบัน ----
 \node[anchor=north east, font=\scriptsize, gray!55!black, align=right] at (4.80,40)
 {27 ส.ค.\\52\% · แผน 65\%};
 % fill=white so the label masks the grid line it sits across
 \node[anchor=north west, font=\scriptsize, kugreen, align=left,
 fill=white, inner sep=1.5pt] at (6.10,68)
 {4 ก.ย. (วันรายงาน)\\ผลจริง 70\% · แผน 72\%};

 % ---- คำอธิบายสัญลักษณ์ (บังกริดด้วยพื้นขาวก่อนวาด) ----
 \fill[white] (5.85,4) rectangle (10.2,26);
 \draw[thick, kudark, densely dashed] (6.0,19) -- (6.7,19);
 \node[right, font=\scriptsize] at (6.8,19) {แผน (Planned)};
 \draw[very thick, kugreen] (6.0,9) -- (6.7,9);
 \node[right, font=\scriptsize] at (6.8,9) {ผลจริง (Actual)};
\end{tikzpicture}
\caption{กราฟ เปรียบเทียบความก้าวหน้าตามแผนกับผลการดำเนินงานจริง
 แกนนอนวางตามขอบ Sprint ทั้ง 6 ช่วง (เส้นประแนวตั้ง)
 จุดบนเส้นผลจริงคือวันที่วัดความก้าวหน้า
 วันรายงาน 4 กันยายน 2569 คือวันที่ 8 จาก 14 ของ Sprint 3 ค่าแผนจึงสอดแทรกจากขอบทั้งสองข้าง
 ช่วงราบระหว่าง 13 ถึง 23 สิงหาคม คือช่วงสอบกลางภาคที่ไม่มีงานเพิ่ม
 ช่วงที่ชันที่สุดคือรอบรายงานนี้}
\label{fig:scurve}
\end{figure}

% =============================================================================
\section{ผลการดำเนินงานเทียบแผนรายวันที่ตั้งไว้}
\label{sec:done}

รายงานครั้งที่ 1 กำหนดงานรายวันไว้ 9 รายการสำหรับสัปดาห์นี้ ตารางนี้รายงานผลของทั้ง 9 รายการตามลำดับเดิม

\noindent\begin{tabularx}{\dimexpr\textwidth-\leftskip\relax}{@{} L{1.3cm} Y L{1.9cm} L{3.2cm} @{}}
\toprule
\hrow \thead{กำหนด} & \thead{งานที่ตั้งไว้} & \thead{ผู้รับผิดชอบ} & \thead{ผล} \\
\midrule
28 ส.ค. & รวมงานฝั่ง backend ที่เสร็จแล้วเข้า main พร้อมตัดสินให้เหลือวิธีจัดระดับอุปกรณ์วิธีเดียว & BE-lead & เสร็จ 30 ส.ค. – 1 ก.ย. ช้ากว่ากำหนด 2–4 วัน เลือกใช้ระดับที่อ่านจากกฎการยืมเป็นระบบเดียว \\
28 ส.ค. & รวมหน้ารับของและหน้าใช้ห้องเข้า main & FE-lead & เสร็จ 30 ส.ค. ช้ากว่ากำหนด 2 วัน \\
29 ส.ค. & Procedure ฝั่งผู้ยืมสำหรับสร้างคำขอยืม ดูคำขอของตนเอง และยกเลิกคำขอ & BE-dev & เสร็จ 2 ก.ย. ช้ากว่ากำหนด 4 วัน และได้ approval queue เพิ่มมาในงานเดียวกัน \\
29 ส.ค. & รวม unique constraint บนอีเมลและรหัสผู้ใช้เข้าสคีมา & SA & เสร็จ 30 ส.ค. ช้ากว่ากำหนด 1 วัน ฐานข้อมูลไม่ยอมรับอีเมลซ้ำอีกต่อไป \\
30 ส.ค. & รับไฟล์ contract ไปใช้แทน contract ที่ฝั่ง frontend เขียนขึ้นเอง & FE-lead & ยังไม่ได้ทำ ติดที่ตัวสร้างไฟล์ contract ยังทำงานไม่ผ่าน \\
31 ส.ค. & เปลี่ยนงานตรวจความสอดคล้องของเวอร์ชันให้ขวางการรวมโค้ด และเพิ่มงานเรียก test & QA-lead & ทำได้ครึ่งเดียว CI เรียก test แล้ว ส่วนงานตรวจเวอร์ชันยังไม่ขวางการรวมโค้ด \\
1 ก.ย. & ติดตั้ง test runner ฝั่ง frontend & QA-lead & ยังไม่ได้ทำ มีเพียงชุด AC check 44 ข้อที่ต้องสั่งรันเอง \\
2 ก.ย. & ต่อหน้ารายการอุปกรณ์และหน้ารายละเอียดอุปกรณ์เข้ากับข้อมูลจริง & FE-lead & เสร็จตามกำหนด \\
3 ก.ย. & ทดสอบวงจรยืม–คืนครบเส้นทางปกติบน main & PM & ยังไม่ได้ทำ สคริปต์ทดสอบที่มีอยู่ครอบคลุมเฉพาะการตั้งค่าคลังอุปกรณ์ ไม่ได้แตะขั้นตอนการยืม \\
\bottomrule
\end{tabularx}

สรุปผลของทั้ง 9 รายการคือเสร็จตามที่ตั้งไว้ 5 รายการ ซึ่งในจำนวนนี้ 4 รายการช้ากว่ากำหนด 1 ถึง 4 วัน ทำได้ครึ่งเดียว 1 รายการ และยังไม่ได้ทำ 3 รายการ

สามรายการที่ยังไม่ได้ทำคือการรับไฟล์ contract การติดตั้ง test runner ฝั่ง frontend และการทดสอบวงจรยืม–คืนครบเส้นทาง ทั้งสามอยู่ใน domain คุณภาพและงาน contract ไม่ใช่งานความสามารถของระบบ ผลคือความสามารถเดินหน้าเร็วกว่าหลักฐานที่ยืนยันว่ามันทำงานถูก ซึ่งเป็นความเสี่ยงที่คณะทำงานบันทึกไว้และแก้ในช่วงที่เหลือของ Sprint ตามหัวข้อ \ref{sec:next}

นอกจากงานตามแผน คณะทำงานยังดึงงานหน้าจอของ Sprint 3 ช่วงหลังขึ้นมาทำก่อนในสองวันสุดท้ายของรอบ ทำให้หน้าจอฝั่งเจ้าหน้าที่และฝั่งผู้อนุมัติทำงานกับข้อมูลจริงได้เร็วกว่าที่วางไว้

\noindent\begin{tabularx}{\dimexpr\textwidth-\leftskip\relax}{@{} L{1.3cm} Y L{1.9cm} L{1.4cm} @{}}
\toprule
\hrow \thead{วันที่} & \thead{งานที่ทำเพิ่มนอกแผน} & \thead{ผู้รับผิดชอบ} & \thead{ผล} \\
\midrule
3–4 ก.ย. & หน้าจอเจ้าหน้าที่ 4 หน้า คือคิวงาน คลังอุปกรณ์ ตรวจสภาพ และค่าตั้งของระบบการยืม ต่อเข้ากับข้อมูลจริงทั้งหมด & FE-lead & เสร็จ \\
4 ก.ย. & หน้าคิวพิจารณาอนุมัติของผู้อนุมัติ ต่อเข้ากับคิวงานและการตัดสินจริง & FE-lead & เสร็จ \\
3 ก.ย. & บังคับ transaction isolation level ให้เข้มขึ้นสำหรับการตรวจการจองชนกัน เพื่อกันสองคำขอของชิ้นเดียวกันที่เข้ามาพร้อมกัน & BE-dev & เสร็จ \\
\bottomrule
\end{tabularx}

% =============================================================================
\section{ผลงานที่ส่งมอบในรอบนี้}

หัวข้อนี้รายงานเฉพาะสิ่งที่ทำงานได้จริงและตรวจสอบได้ใน repository ไม่รวมกิจกรรมระหว่างดำเนินงาน

\subsection{ฝั่ง backend}

\begin{itemize}
 \item วงจรยืม–คืนมี procedure ครบทุกขั้น — ผู้ยืมสร้างคำขอและยกเลิกคำขอของตนเองได้ · ผู้อนุมัติเห็นคิวงานพร้อมจำนวนคงค้างและตัดสินได้ · เจ้าหน้าที่จัดของ ผูกหมายเลขชิ้น สลับชิ้นที่จ่าย ยืนยันการรับของ บันทึกการคืน และบันทึกกรณีสูญหาย · ตรวจสภาพพร้อมแนบรูป ส่งซ่อม และเสนอปลดระวาง
 \item กฎธุรกิจรายระดับ — อุปกรณ์และห้องถูกจัดระดับ T0 ถึง T3 จากกฎการยืมเพียงระบบเดียว ปิดข้อขัดแย้งที่รอบก่อนระบุว่ามีวิธีจัดระดับสองวิธีวางซ้อนกัน · เส้นทางอนุมัติและจำนวนวันยืมสูงสุดคำนวณจากระดับนี้
 \item การจองสถานที่ (T3) — ห้องใช้เส้นทางเดียวกับอุปกรณ์ทั้งหมด โดยมี buffer ก่อนและหลังการใช้งานตามที่กำหนดไว้กับทรัพยากรแต่ละรายการ และกันการจองซ้อนช่วงเวลาโดยนับคำขอที่ยังรอพิจารณาว่าถือครองอยู่ด้วย
 \item ระบบ notification — ผูกกับเหตุการณ์จริงในวงจรยืมและการอนุมัติ ไม่ใช่ mock · มีกุญแจกันการแจ้งซ้ำ ทำให้เหตุการณ์เดียวกันไม่ถูกแจ้งสองครั้ง · อ่านรายการ นับจำนวนที่ยังไม่อ่าน และทำเครื่องหมายว่าอ่านแล้วได้
 \item ความปลอดภัยของบัญชีผู้ใช้ — ตาราง session ทำให้เพิกถอนสิทธิ์ได้ทันทีและออกจากระบบทุกอุปกรณ์พร้อมกันได้ · จำกัดจำนวนครั้งการเข้าสู่ระบบที่ผิดพลาด · audit log ทุกการกระทำที่สำคัญ · จำกัดขอบเขตข้อมูลที่เจ้าหน้าที่แต่ละคนเห็นตามหน่วยงานที่ดูแล · unique constraint บนอีเมลและรหัสผู้ใช้
 \item ข้อมูล seed ครบชุด — อุปกรณ์ ห้อง กฎการยืม ระดับเครดิต และบัญชีครบทุกบทบาท พร้อมสคริปต์เดินวงจรยืม–คืนอัตโนมัติเพื่อยืนยันว่าข้อมูลชุดนี้ใช้ทำงานได้จริง
 \item ขนาดของสิ่งที่ส่งมอบ ณ วันรายงาน — 9 domain รวม 78 procedure บนฐานข้อมูล 29 ตาราง 9 Enum
\end{itemize}

ข้อจำกัดที่เหลือมีห้าข้อ ข้อแรกและสำคัญที่สุดคือ procedure ของวงจรยืม–คืน การอนุมัติ และ notification ยังไม่มี test แม้แต่กรณีเดียว สิ่งที่ยืนยันได้ ณ วันรายงานคือ procedure เหล่านี้ build และ mount route ได้เท่านั้น ข้อที่สองคือกลุ่มการอุทธรณ์และกลุ่มรายงานเชิงสถิติยังไม่มีฝั่ง backend เลย ข้อที่สามคือค่าตั้งเวลาทำการยังไม่มีที่เก็บในฐานข้อมูล procedure อ่านและแก้ค่าจึงยังตอบกลับว่ายังไม่รองรับ ข้อที่สี่คือ cron job ยังไม่ทำงานเอง ต้องสั่งด้วยมือ และข้อที่ห้าคือไฟล์รูปที่แนบตอนตรวจสภาพถูกเก็บบนดิสก์ของเครื่องและเปิดอ่านได้โดยไม่ตรวจสิทธิ์ ซึ่งยังไม่ตรงกับที่ SRS กำหนดไว้ว่าให้เก็บบนที่เก็บภายนอกและเข้าถึงผ่าน pre-signed URL

\subsection{ฝั่ง frontend}

\begin{itemize}
 \item หน้าจอที่มีเนื้อหาจริงเพิ่มเป็น 24 จาก 30 หน้า — เพิ่มหน้าแรกของผู้ยืม หน้ารับของ หน้าใช้ห้อง และหน้าอุทธรณ์ พร้อมปุ่มขอต่อเวลาสำหรับอุปกรณ์ระดับ T1 และ T2 · และเพิ่มฝั่งเจ้าหน้าที่กับฝั่งผู้อนุมัติในช่วงท้ายรอบ คือคิวงานของเจ้าหน้าที่ คลังอุปกรณ์ ตรวจสภาพ ค่าตั้งของระบบการยืม และคิวพิจารณาอนุมัติ โดยยุบหน้าแผงสรุปของเจ้าหน้าที่เข้ากับหน้าคิวงาน เพราะทั้งสองหน้าแสดงข้อมูลชุดเดียวกัน
 \item หน้าจอที่อ่านข้อมูลจริงเพิ่มเป็น 13 หน้า — เข้าสู่ระบบ · รายการอุปกรณ์ · รายละเอียดอุปกรณ์ · สร้างคำขอ · โปรไฟล์และคะแนนเครดิต · หน้าแรกของผู้ยืม · จัดการผู้ใช้ · audit log · คิวงานเจ้าหน้าที่ · คลังอุปกรณ์ · ตรวจสภาพ · ค่าตั้งของระบบการยืม · คิวพิจารณาอนุมัติ รอบก่อนมีเพียงเส้นทางเข้าสู่ระบบเส้นทางเดียว
 \item กระดิ่ง notification บน sidebar อ่านจากระบบ notification จริงและใช้ได้ทุกหน้า
 \item ตัวช่วยไล่อ่านข้อมูลทีละหน้า (pagination) สำหรับรายการที่ฝั่ง backend จำกัดจำนวนแถวต่อครั้ง
 \item ชุด AC check ฝั่ง frontend 44 ข้อ ซึ่งสั่งรันแล้วผ่านทั้งหมด ครอบคลุมช่วงเวลายืม เวลาตัดคืน 17:00 น. การนับช่องจองพร้อมกันของ T3 การตรวจไฟล์อัปโหลด และการแปลข้อความแสดงข้อผิดพลาด
\end{itemize}

ข้อจำกัดที่เหลือมีสี่ข้อ ข้อแรกคือปุ่มส่งคำขอของผู้ยืมยังบันทึกคำขอลง store ในเบราว์เซอร์แทนที่จะส่งไปยังฐานข้อมูล คำขอจึงหายเมื่อรีเฟรชหน้า และเป็นจุดเดียวที่ทำให้วงจรยืม–คืนยังเดินครบผ่านหน้าจอไม่ได้ ทั้งที่ปลายทางฝั่งเจ้าหน้าที่และผู้อนุมัติพร้อมแล้ว ข้อที่สองคือหน้าจอที่มี UI แล้วอีก 11 หน้ายังอ่าน mock และยังเหลือหน้าว่างอีก 6 หน้า คือเจ้าหน้าที่ 3 หน้า (ส่งมอบ สิทธิ์ ผู้ใช้) รายงาน 2 หน้า และตั้งค่าผู้ดูแลระบบ 1 หน้า ข้อที่สามคือฝั่ง frontend ยังใช้ไฟล์ contract ที่เขียนขึ้นเอง ซึ่งประกาศ domain ที่ฝั่ง backend ไม่มีอยู่ 3 กลุ่มและตั้งชื่อ procedure ไม่ตรงกับของจริงอีกหลายจุด รวมแล้วมี 15 procedure ที่ frontend ประกาศไว้แต่เรียกใช้จริงไม่ได้ ผลคือโค้ดผ่านการ typecheck ทั้งที่เรียกสิ่งที่ไม่มีอยู่ และจะพังตอนผู้ใช้กดปุ่มเท่านั้น ข้อที่สี่คือยังไม่มี test runner ชุดตรวจ 44 ข้อจึงต้องสั่งรันเอง

\subsection{เครื่องมือและกระบวนการ}

CI เดินครบห้าขั้นตอนแล้ว คือติดตั้ง typecheck build migrate และเรียก test 227 กรณีบนฐานข้อมูล PostgreSQL จริงที่ยกขึ้นมาเฉพาะการตรวจแต่ละครั้ง ข้อจำกัดที่รอบก่อนระบุไว้ว่า ``test พังได้โดยไม่มีใครรู้'' จึงปิดแล้วสำหรับฝั่ง backend

ข้อจำกัดที่เหลือมีสองข้อ ข้อแรกคืองานตรวจความสอดคล้องของเวอร์ชันไลบรารีร่วมยังตั้งไว้ให้รายงานผลโดยไม่ขวางการรวมโค้ด ทั้งที่เงื่อนไขที่เคยต้องผ่อนผันหมดไปแล้วตั้งแต่รอบก่อน ข้อที่สองคือการตรวจฝั่ง frontend ยังมีแค่ typecheck กับสร้างไฟล์ ยังไม่มีขั้นตอนเรียก test

\subsection{เอกสารและการออกแบบ}

\begin{itemize}
 \item เอกสารข้อกำหนดความต้องการของซอฟต์แวร์ (SRS) ฉบับทางการ — เรียงตามกรอบ IEEE 830-1998 ครอบคลุมรหัสความต้องการของทีมครบทุกข้อ ส่งเป็นไฟล์ที่แก้ไขต่อได้พร้อมฉบับ PDF จากต้นฉบับเดียวกัน
 \item เอกสารการออกแบบระบบ (SDS) ฉบับทางการ — เขียนจากโค้ดจริงบน main ไม่ใช่จากสิ่งที่คาดว่าจะเป็น สิ่งที่ยังไม่ได้ทำระบุไว้ในคอลัมน์สถานะและภาคผนวก เรียงตามโครงที่วิชากำหนดครบทั้ง 11 หัวข้อ
 \item เอกสารอ้างอิง API ที่สร้างจากโค้ดโดยตรง จึงไม่มีวันคลาดจากสิ่งที่ระบบทำจริง พร้อมบันทึกการตัดสินใจเรื่องวิธียืนยันตัวตน และเอกสารอธิบายงานฝั่งเจ้าหน้าที่กับฝั่งผู้ดูแลระบบ
\end{itemize}

ข้อจำกัดที่เหลือคือยังไม่มีคู่มือผู้ใช้และคู่มือติดตั้ง ทั้งสองรายการอยู่ในแผน Sprint 5 และ SRS กับ SDS ต้องปรับตามเมื่อกลุ่มการอุทธรณ์และกลุ่มรายงานถูกเพิ่มเข้าระบบ

\subsection{ระดับหลักฐานของสิ่งที่รายงานว่าเสร็จ}

หัวข้อนี้เพิ่มขึ้นในรอบนี้เพื่อไม่ให้คำว่า ``เสร็จ'' ถูกอ่านเกินกว่าที่หลักฐานรองรับ คอลัมน์สุดท้ายแยกสองอย่างออกจากกัน คือเขียนโค้ดเสร็จ กับมีอะไรพิสูจน์แล้วว่าทำงานถูก

\noindent\begin{tabularx}{\dimexpr\textwidth-\leftskip\relax}{@{} L{3.8cm} Y L{3.0cm} @{}}
\toprule
\hrow \thead{สิ่งที่ส่งมอบ} & \thead{หลักฐาน ณ วันรายงาน} & \thead{ระดับ} \\
\midrule
contract 78 procedure ใน 9 domain & ระบบ mount route ครบทุก procedure ตอนเริ่มทำงาน และ CI ผ่าน & พิสูจน์แล้ว \\
ยืนยันตัวตน session และ audit log & test 26 กรณีทำงานบนฐานข้อมูลจริงใน CI & พิสูจน์แล้ว \\
รายการอุปกรณ์ ห้อง และการจำกัดสิทธิ์ข้ามหน่วยงาน & สคริปต์ 8 ขั้นตอนบนฐานข้อมูลจริง ต้องสั่งรันเอง ไม่ได้อยู่ใน CI & พิสูจน์แล้ว แต่ไม่อัตโนมัติ \\
กฎธุรกิจฝั่ง frontend & ชุด AC check 44 ข้อ ผ่านทั้งหมด ต้องสั่งรันเอง & พิสูจน์แล้ว แต่ไม่อัตโนมัติ \\
วงจรยืม–คืน การอนุมัติ และ notification & ไม่มี test และไม่มีสคริปต์ใดเรียกใช้ ตรวจได้เพียงว่าคอมไพล์และ build ผ่าน & เขียนเสร็จ ยังไม่พิสูจน์ \\
หน้าจอเจ้าหน้าที่และผู้อนุมัติ & ต่อกับ procedure จริงครบทั้งการอ่านและการสั่งงาน typecheck ผ่าน แต่ยังไม่มีการทดสอบหน้าจอ & เขียนเสร็จ ยังไม่พิสูจน์ \\
\bottomrule
\end{tabularx}

สองแถวสุดท้ายคือส่วนที่เขียนขึ้นในรอบนี้เอง และเป็นส่วนที่มีหลักฐานน้อยที่สุด คณะทำงานถือว่าเป็นความเสี่ยงหลักที่ยกไปแก้ในช่วงที่เหลือของ Sprint 3 ไม่ใช่สิ่งที่นับว่าเสร็จแล้ว

% =============================================================================
\section{แผนการดำเนินงานถัดไป}
\label{sec:next}

\subsection{แผนช่วงที่เหลือของ Sprint 3 (4 – 10 กันยายน 2569)}

ช่วงที่เหลือของ Sprint จบลงที่วันสิ้นสุด Sprint พอดี งานเกือบทั้งหมดเป็นการต่อหน้าจอเข้ากับความสามารถที่มีอยู่แล้ว ไม่ใช่การเพิ่มความสามารถใหม่ ยกเว้นกลุ่มการอุทธรณ์

งานสองรายการแรกแก้ที่ต้นเหตุของรายการที่ค้างมาสองรอบ คือฝั่ง frontend ยังไม่ได้รับไฟล์ contract ไปใช้ ระหว่างตรวจสอบเพื่อทำรายงานฉบับนี้ คณะทำงานได้สั่งรันตัวสร้าง contract ดูจริง และพบว่ามันทำงานได้ปกติแล้ว สร้างไฟล์ที่ประกาศ 9 domain 78 procedure ครบถ้วน อุปสรรคที่เคยบันทึกไว้ว่าตัวสร้างติดคำอธิบายภาษาไทยนั้นเกิดเฉพาะเมื่อไฟล์ router ประกาศ schema ไว้ในตัวเอง ซึ่งขณะนี้ไม่มี router ไฟล์ใดทำเช่นนั้นแล้ว ทุกไฟล์ย้ายไปอ้างถึง schema จากภายนอกทั้งหมด เงื่อนไขที่ทำให้ตัวสร้างล้มจึงหมดไปเอง

อุปสรรคที่แท้จริงคือไฟล์ที่สร้างออกมาอ้างถึง schema ในโค้ด backend ด้วยพาธสัมพัทธ์ ฝั่ง frontend ซึ่งเป็นคนละโปรเจกต์จึงยังตามพาธเหล่านั้นไม่เจอ งานที่เหลือคือทำให้ฝั่ง frontend เข้าถึง schema ชุดนั้นได้ ตามแนวทางที่ทีมเลือกไว้ตั้งแต่ Sprint 1 คือแยก package แล้วล็อกเวอร์ชัน ไม่ใช่การไล่แก้คำอธิบายในโค้ด

\noindent\begin{tabularx}{\dimexpr\textwidth-\leftskip\relax}{@{} L{1.5cm} Y L{2.1cm} @{}}
\toprule
\hrow \thead{กำหนด} & \thead{งาน} & \thead{ผู้รับผิดชอบ} \\
\midrule
4 ก.ย. & สร้างไฟล์ contract จากโค้ดจริงเข้า main และทำให้ฝั่ง frontend อ้างถึง schema ที่ไฟล์นั้นต้องใช้ได้ ซึ่งเป็นงานที่เหลืออยู่จริงหลังพบว่าตัวสร้างทำงานได้อยู่แล้ว & BE-lead \\
5 ก.ย. & รับไฟล์ contract ไปใช้แทนไฟล์ที่ฝั่ง frontend เขียนขึ้นเอง และลบไฟล์เดิมทิ้ง เพื่อให้การเรียกสิ่งที่ฝั่ง backend ไม่มีถูกจับได้ตั้งแต่ตอนคอมไพล์ & FE-lead \\
6 ก.ย. & ต่อปุ่มส่งคำขอของผู้ยืมเข้ากับ procedure สร้างคำขอจริง และเลิกใช้ store ในเบราว์เซอร์ ซึ่งเป็นจุดเดียวที่ยังขวางไม่ให้เดินวงจรครบผ่านหน้าจอ & FE-dev \\
7 ก.ย. & หน้าส่งมอบของเจ้าหน้าที่ ซึ่งเป็นหน้าว่างหน้าสุดท้ายที่อยู่บนเส้นทางของวงจรยืม–คืน & FE-lead \\
8 ก.ย. & test ที่เดินวงจรยืม–คืนตั้งแต่สร้างคำขอจนรับคืนบนฐานข้อมูลจริง เพื่อให้ส่วนที่เขียนใหม่มีหลักฐานเท่ากับส่วนที่เขียนไว้ก่อน & BE-lead \\
8 ก.ย. & ติดตั้ง test runner ฝั่ง frontend ย้ายชุด AC check 44 ข้อเข้า CI และเปลี่ยนงานตรวจเวอร์ชันให้ขวางการรวมโค้ด & QA-lead \\
9 ก.ย. & ต่อหน้ารายการยืมของฉัน หน้ารายการห้อง และหน้ารับของ เข้ากับข้อมูลจริง & FE-lead \\
9 ก.ย. & กลุ่มการอุทธรณ์ฝั่ง backend และที่เก็บค่าตั้งเวลาทำการ & BE-dev \\
10 ก.ย. & เดินวงจรยืม–คืนครบเส้นทางปกติผ่านหน้าจอจริง ตั้งแต่ผู้ยืมสร้างคำขอจนเจ้าหน้าที่รับคืน ซึ่งเป็นเกณฑ์ ``เสร็จ'' ของ Sprint 3 & PM \\
\bottomrule
\end{tabularx}

\subsection{Sprint 4 และ Sprint 5}

\noindent\begin{tabularx}{\dimexpr\textwidth-\leftskip\relax}{@{} L{2.9cm} Y @{}}
\toprule
\hrow \thead{Sprint} & \thead{เป้าหมายและเกณฑ์ ``เสร็จ''} \\
\midrule
\textbf{Sprint 4} \newline 11 – 24 ก.ย. & ข้ามหน่วยงานและงานขัดเงา — ตะกร้าพร้อมการแตกอนุมัติตามหน่วยงานเจ้าของ · การจองล่วงหน้าถึงสิ้นภาคการศึกษาสำหรับ T1–T3 · กลุ่มรายงานเชิงสถิติและหน้าจอรายงาน 2 หน้า · หน้าจอที่ยังว่างที่เหลือ · การค้นหา การแบ่งหน้า และการรองรับมือถือ \newline เกณฑ์เสร็จ — ไม่เหลือหน้าว่างในเส้นทางที่ผู้ใช้จริงต้องใช้ และไม่เหลือ mock ในระบบ \newline หมายเหตุ — เป็น Sprint สุดท้ายที่เพิ่มความสามารถใหม่ได้ \\
\textbf{Sprint 5} \newline 25 ก.ย. – 1 ต.ค. & หยุดเพิ่มความสามารถและเตรียมนำเสนอ — ทดสอบถดถอยเต็มรูปแบบ · ทดสอบสมรรถนะ · นำระบบขึ้นใช้งานจริงพร้อมย้ายที่เก็บไฟล์รูปไปไว้ภายนอกตามที่ SRS กำหนด · ปรับ SRS และ SDS ให้ตรงกับระบบที่ส่งมอบ · คู่มือผู้ใช้และคู่มือติดตั้ง · ซ้อมนำเสนอ \\
\bottomrule
\end{tabularx}

\subsection{การประเมินกำหนดเสร็จ}

จุดวัดสามข้อที่ตั้งไว้ในรายงานครั้งที่ 1 ข้อแรกผ่านแบบช้ากว่ากำหนด 1 ถึง 3 วัน คืองานที่เสร็จแล้วถูกรวมเข้า main ครบทุกก้อน ข้อที่สองไม่ผ่าน คือวงจรยืม–คืนยังไม่ได้ถูกทดสอบครบเส้นทาง ทั้งผ่านหน้าจอและด้วย test ส่วนข้อที่สามยังไม่ถึงกำหนด

คณะทำงานประเมินว่าโครงการยังเสร็จได้ภายใน 1 ตุลาคม 2569 ตามกำหนดเดิม โดยมีขอบเขตครบตามข้อเสนอโครงการ เหตุผลคือความเสี่ยงเปลี่ยนลักษณะไปจากรอบก่อน รอบก่อนความเสี่ยงอยู่ที่ความสามารถของระบบซึ่งยังไม่มี ส่วนรอบนี้ความสามารถหลักเขียนครบแล้ว และความเสี่ยงย้ายไปอยู่ที่การต่อหน้าจอกับการยืนยันว่าสิ่งที่เขียนไว้ทำงานถูก ซึ่งเป็นงานที่ประเมินเวลาได้แม่นกว่าและทำคู่ขนานกันได้หลายคน

ข้อควรระวังที่คณะทำงานบันทึกไว้มีสองข้อ ข้อแรกคือความชันของรอบนี้ส่วนหนึ่งมาจากงานที่เขียนเสร็จก่อนหน้าแล้วเพิ่งถูกนับเมื่อรวมเข้า main จึงไม่ควรใช้อัตราของรอบนี้ประมาณการรอบถัดไป ข้อที่สองคือความสามารถของระบบเดินหน้าเร็วกว่าหลักฐานที่ยืนยันว่ามันทำงานถูก ส่วนที่เขียนขึ้นใหม่ในสองสัปดาห์หลังยังไม่มี test ตัวเลขความก้าวหน้าในรายงานนี้จึงสะท้อนปริมาณงานที่เขียนเสร็จ ไม่ใช่ปริมาณงานที่ตรวจสอบแล้วว่าถูกต้อง การประเมินข้างต้นตั้งอยู่บนจุดวัดสี่ข้อถัดไป ถ้าข้อใดไม่ผ่าน คณะทำงานจะทบทวนลำดับความสำคัญของ Sprint 4 ทันทีในการทบทวนปิด Sprint 3

\begin{enumerate}
 \item ปุ่มส่งคำขอของผู้ยืมเขียนคำขอลงฐานข้อมูลจริง และหน้าส่งมอบของเจ้าหน้าที่ใช้งานได้ ภายใน 7 กันยายน 2569
 \item เดินวงจรยืม–คืนครบเส้นทางปกติผ่านหน้าจอจริงได้ ภายในวันสิ้นสุด Sprint 3 คือ 10 กันยายน 2569
 \item วงจรยืม–คืนมี test อัตโนมัติที่เดินครบเส้นทางบนฐานข้อมูลจริง ภายในวันสิ้นสุด Sprint 3
 \item ไม่เหลือ mock ในเส้นทางหลัก และฝั่ง frontend ใช้ไฟล์ contract ที่สร้างจากโค้ดจริง ภายใน 14 กันยายน 2569
\end{enumerate}

% =============================================================================
\section*{ภาคผนวก — ค่าที่ใช้พล็อตกราฟ}

ค่าเป้าหมายรายมิติ ณ วันสิ้นสุดของแต่ละ Sprint คูณด้วยน้ำหนักของมิตินั้นแล้วรวมกัน คอลัมน์ ``รวม'' คือค่าที่พล็อตบนเส้นแผนในรูปที่ \ref{fig:scurve}

แถวที่มีเครื่องหมายดอกจันคือวันรายงานรอบนี้ ซึ่งอยู่ที่ 8 ใน 14 วันของ Sprint 3 ค่าในแถวนี้จึงเป็นค่าสอดแทรกเชิงเส้นของสองแถวที่ขนาบอยู่ตามสัดส่วนนั้น ปัดเป็นจำนวนเต็มรายมิติก่อนถ่วงน้ำหนัก

แถวแรก ``เริ่ม 24 ก.ค.'' คือจุดตั้งต้นของการนับ ทุกมิติเป็นศูนย์ยกเว้น Documentation ที่ 10\% ซึ่งมาจากข้อเสนอโครงการ งานออกแบบที่เสร็จก่อนวันนั้นถือเป็นข้อมูลนำเข้าของงานออกแบบใน Sprint 0 ไม่ใช่ค่าตั้งต้น

\noindent\begin{tabularx}{\dimexpr\textwidth-\leftskip\relax}{@{} L{3.2cm} L{1.0cm} L{1.0cm} L{1.0cm} L{1.0cm} L{1.0cm} L{1.0cm} Y @{}}
\toprule
\hrow \thead{ขอบ Sprint} & \thead{BE Des} & \thead{FE Des} & \thead{BE Dev} & \thead{FE Dev} & \thead{Test} & \thead{Doc} & \thead{รวม} \\
\hrow \thead{} & \thead{10\%} & \thead{10\%} & \thead{30\%} & \thead{25\%} & \thead{15\%} & \thead{10\%} & \thead{100\%} \\
\midrule
เริ่ม 24 ก.ค. & 0 & 0 & 0 & 0 & 0 & 10 & 1 \\
S0 · 30 ก.ค. & 45 & 30 & 5 & 5 & 0 & 20 & 12 \\
S1 · 13 ส.ค. & 88 & 75 & 25 & 42 & 10 & 65 & 42 \\
S2 · 27 ส.ค. & 95 & 85 & 60 & 65 & 35 & 70 & 65 \\
$*$ วันรายงาน 4 ก.ย. & 98 & 89 & 70 & 72 & 45 & 73 & 72 \\
S3 · 10 ก.ย. & 100 & 92 & 78 & 78 & 52 & 76 & 78 \\
S4 · 24 ก.ย. & 100 & 100 & 97 & 97 & 80 & 88 & 94 \\
S5 · 1 ต.ค. & 100 & 100 & 100 & 100 & 100 & 100 & 100 \\
\bottomrule
\end{tabularx}

ค่าผลจริง ณ วันรายงานที่ใช้พล็อตเส้นผลจริงในรูปที่ \ref{fig:scurve} และแท่งในรูปที่ \ref{fig:dims} คือ 95, 91, 69, 69, 32 และ 85 ตามลำดับมิติเดียวกัน ถ่วงน้ำหนักแล้วได้ \actualpct

\end{document}
